%==============================================================================
% Chapitre 70 - La mise en mouvement du projectile
%==============================================================================

\section{Introduction}

Pendant les premiers instants du tir, le projectile demeure immobile malgré la
combustion de la poudre et l'augmentation rapide de la pression.

Le mouvement ne débute que lorsque les efforts exercés par les gaz deviennent
supérieurs aux différentes résistances qui s'opposent au déplacement.

\begin{aretenir}

Le projectile ne commence pas à se déplacer dès l'inflammation de la poudre.
Il reste immobile jusqu'à ce que la force résultante devienne suffisante.

\end{aretenir}

\section{Application de la pression}

La pression des gaz s'exerce sur toute la surface du culot du projectile.

Cette pression engendre une force dirigée vers l'avant, dans l'axe du canon.

Plus la pression augmente, plus cette force devient importante.

\section{Les résistances au déplacement}

Avant que le projectile ne puisse avancer, plusieurs résistances doivent être
vaincues :

\begin{itemize}
    \item son inertie ;
    \item le sertissage éventuel ;
    \item les frottements avec le collet de l'étui ;
    \item l'engagement dans les rayures du canon.
\end{itemize}

Ces résistances expliquent pourquoi la pression continue d'augmenter pendant les
premiers instants de la combustion.

\section{Le rôle de l'inertie}

Selon le principe d'inertie, un corps immobile tend à conserver son état de
repos tant qu'une force suffisante ne s'exerce pas sur lui.

Le projectile ne fait pas exception.

Sa masse influence directement l'effort nécessaire pour initier son mouvement.

\section{Gravure dans les rayures}

Dans une arme rayée, le projectile doit s'engager dans les rayures du canon.

Cette opération nécessite un effort supplémentaire.

Elle constitue l'une des principales résistances rencontrées au début du
déplacement.

\section{Début du mouvement}

Lorsque la force exercée par les gaz dépasse l'ensemble des résistances, le
projectile commence à avancer.

Ce déplacement reste d'abord très faible, mais il augmente rapidement sous
l'effet de la pression.

\section{Conséquence sur le volume}

Le mouvement du projectile augmente progressivement le volume disponible pour
les gaz.

Cette augmentation modifie l'évolution de la pression dans la chambre.

La combustion et le déplacement du projectile deviennent alors fortement
couplés.

\section{Équilibre dynamique}

Pendant toute cette phase, deux phénomènes évoluent simultanément :

\begin{itemize}
    \item la combustion continue à produire des gaz ;
    \item le déplacement du projectile augmente le volume disponible.
\end{itemize}

L'évolution de la pression résulte de cet équilibre dynamique.

\section{Influence de la masse du projectile}

À charge propulsive identique, un projectile plus lourd nécessite généralement
une pression plus importante pour atteindre une même accélération initiale.

La masse du projectile influence donc directement le comportement de la
munition.

\section{Influence du sertissage}

Le sertissage contribue à maintenir le projectile dans l'étui avant le tir.

Il augmente légèrement la résistance initiale au déplacement.

Son influence est prise en compte lors de la conception des munitions afin
d'assurer une combustion régulière.

\section{Observation expérimentale}

Les premières phases du déplacement peuvent être étudiées à l'aide de
caméras ultra-rapides, de capteurs de pression et de mesures de vitesse.

Ces observations permettent de valider les modèles de balistique intérieure.

\section{Représentation en simulation}

\begin{simulation}

Dans les modèles numériques, le projectile est généralement considéré comme
immobile tant que la force exercée par les gaz reste inférieure aux résistances
initiales.

Lorsque ce seuil est dépassé, le calcul de son mouvement débute à chaque pas de
temps selon les lois de la mécanique.

\end{simulation}

\section{Vers l'accélération}

Une fois le projectile en mouvement, les gaz continuent à exercer une force sur
son culot.

Il entre alors dans une phase d'accélération rapide qui se poursuivra jusqu'à
sa sortie du canon.

\section{À retenir}

\begin{aretenir}

\begin{itemize}
    \item Le projectile reste immobile durant les premiers instants du tir.
    \item La pression agit sur son culot.
    \item Plusieurs résistances doivent être vaincues avant le déplacement.
    \item Le mouvement du projectile augmente progressivement le volume des gaz.
    \item La combustion et le déplacement évoluent simultanément.
\end{itemize}

\end{aretenir}

\begin{rigueur}{Pression et force}

La pression est une grandeur répartie sur une surface.

La force exercée sur le projectile dépend de la pression ainsi que de la surface
du culot.

Cette distinction est essentielle pour comprendre les équations introduites
dans les chapitres suivants.

\end{rigueur}

\begin{ideerecue}

Un projectile ne commence pas immédiatement à se déplacer lorsque l'amorce
fonctionne.

\end{ideerecue}

\begin{ordredegrandeur}

Le projectile reste immobile pendant quelques dizaines à quelques centaines de
microsecondes avant de commencer son déplacement.

\end{ordredegrandeur}

\begin{ordredegrandeur}

La pression maximale atteint couramment plusieurs centaines de mégapascals.

\end{ordredegrandeur}
